% experiments/exp2_chain_recipe/intro_v3.tex
%
% Replacement Section 1 (Introduction) for draft.tex, restructured
% around the question "When is recurrent memory necessary, and how
% do we know?" and incorporating the v14/v15 findings of
% section7_v14_findings.tex, section8_recipe_revised.tex, and
% section_eval_callback.tex.
%
% This file is meant to be \input'd in place of the original
% \section{Introduction} block of draft.tex. It does not modify
% draft.tex.

\section{Introduction}
\label{sec:intro}

A practical conversational agent should remember a user across
sessions: their name, their preferences, the on-going arc of a
relationship. The contemporary toolkit offers three options for
this, and none is good. \emph{Long context} is expensive per turn
and degrades with distance. \emph{Retrieval-augmented generation}
works when the right token is known to live in a known document,
but fails on the recall pattern that matters most in dialogue ---
\emph{abstract callback}, where what must be retrieved is not a
substring of any prior turn but the gist of an evidence cluster
spread across them. \emph{Recurrent memory} should fit the
abstract-callback pattern by construction: a fixed-size state
$\Mc$ updated end-to-end through the language-modelling loss, with
no retrieval index, no separate controller, and no hand-engineered
gating. Whether it does in practice is the question this paper
attacks. The frame we adopt is methodological: \emph{when is
recurrent memory necessary, and how do we know it is helping?}

\paragraph{Where the prior work leaves us.} The companion drop-in
primitive paper~\cite{liu2025drop} establishes that the
\emph{Memory Residuals} primitive, on a $0.6$B Qwen3 backbone
trained on book continuations, is roughly $2\!\times$ more
sample-efficient on the history-specificity diagnostic ($\dsh$,
\S\ref{sec:primitive}) than its ReZero-style gate baseline. That
is an \emph{in-trainer, phase-aligned} signal: $\dnm$ and $\dsh$
are computed under the same data sampler and session boundaries
the optimiser saw. When the same checkpoint is taken through a
standalone eval pipeline on dialogue, the signal becomes mixed.
The first iteration of the present manuscript
(\S\S\ref{sec:recipe}--\ref{sec:experiments}) gives a four-piece
recipe that survives this transfer, plus an analytic negative
result on the strict $\pm 32$ parity init
(\S\ref{sec:negative_result}). The second iteration
(\S\ref{sec:second_iteration}) adds three training-dynamics pieces
--- supervised writer warmup, $\amem$ load-balance auxiliary, and
multi-negative InfoNCE --- each motivated by a diagnostic signature
on a heterogeneous mixed corpus. By the end of the second iteration
the recipe has seven pieces, and the question of whether memory is
helping on dialogue still hinges on which subset of the score
window the eval looks at.

\paragraph{The pivot.} A controlled synthetic callback corpus
(D4v2, \S\ref{sec:third_iteration}) and a callback-aware standalone
eval (\S\ref{sec:eval_cb}) collapse the ambiguity. On D4v2 a
six-cell v14 ablation establishes that the supervised writer
warmup is empirically counterproductive: it imprints a
``compress everything'' policy on the writer that, post-anneal,
locks the router into a recent-bias regime with negative
evidence-lift on the very tokens the warmup was meant to help.
Removing the warmup, with input-RMSNorm on the extractor and the
alpha-floor and InfoNCE pieces in place, yields the first sustained
positive evidence-lift we have observed on a controlled callback
corpus. Standalone evaluation of that checkpoint then surfaces the
methodological finding: \emph{standard} chain eval averages CE
uniformly across the score window, dilutes the callback-positioned
signal across roughly four sessions of filler-and-evidence, and
drives its signal-to-noise below one. A callback-aware eval
recovers a $+1.4$~nat memory-on-callback effect on the same
checkpoint that the full-window eval averages to noise.

\paragraph{Contributions.} This paper extends the four-piece recipe
of \S\ref{sec:recipe} along five axes:
\begin{enumerate}\setlength{\itemsep}{0pt}
\item \textbf{The v14 warmup pathology
(\S\ref{sec:third_iteration}).} A six-cell ablation on D4v2 shows
that the writer warmup of \S\ref{sec:second_iteration} is at best
inert and at worst silently counterproductive, with a mechanical
explanation tied to the writer's response to a saturated router
under a frozen backbone.
\item \textbf{The revised three-piece recipe
(\S\ref{sec:revised_recipe}).} Input-RMSNorm on the extractor, no
writer warmup, alpha-floor $+$ multi-negative InfoNCE: the
minimal recipe under which we observe sustained positive
evidence-lift on D4v2 at $0.6$B frozen.
\item \textbf{Callback-aware standalone eval
(\S\ref{sec:eval_cb}).} A standalone-eval methodology that
restricts CE scoring to the callback session and the annotated
callback-position tokens within it, and that measures the
$\dnm$-floor against a memory state with the chain's evidence
sessions redacted by distractor sessions sampled from filler
positions of other chains. Released as
\texttt{tools/eval\_callback.py}.
\item \textbf{Honest negative on joint training at $0.6$B.} With
the backbone unfrozen, the bare backbone learns the D4v2 template
directly and the memory signal vanishes regardless of recipe; we
document the failure mode and discuss its implications for the
freeze-vs-joint axis (v15b).
\item \textbf{Scaling probe at $1.7$B (pending,
Table~\ref{tab:revised_headline}).} v15e (frozen) and v15f (joint)
port the revised three-piece recipe to Qwen3-1.7B; whether the
freeze-required-for-memory boundary moves with scale is the
headline open question.
\end{enumerate}

\paragraph{Roadmap.} \S\ref{sec:primitive} folds in the primitive's
definition and book-only baseline from \cite{liu2025drop}.
\S\ref{sec:recipe}--\ref{sec:setup} state the original four-piece
recipe and the corpus build. \S\ref{sec:experiments} reports the
$2\!\times\!2$ headline. \S\ref{sec:negative_result} pins down
the strict-vs-soft init failure. \S\ref{sec:second_iteration}
reports the second-iteration three-piece extension.
\S\ref{sec:third_iteration} reports the v14 ablation matrix on
D4v2 and the warmup catastrophe. \S\ref{sec:revised_recipe}
states the revised three-piece recipe and the v15 scaling cells.
\S\ref{sec:eval_cb} introduces the callback-aware standalone eval
and contrasts it numerically with the standard chain eval on the
v14k checkpoint. \S\ref{sec:discussion} closes with limitations
and the interaction of architectural, training-recipe, and
evaluation-methodology interventions.
